\chapter{Steady RANS Validation at Cruise Conditions}
\label{ch:rans_validation}

% =====================================================================
% SKELETON - ogni sezione contiene TODO con cosa scrivere.
% Il paragrafo originale e' conservato in fondo alla sezione 1.
% =====================================================================

\section{Introduction and objectives}
\label{sec:val_intro}

Since its introduction for the fourth AIAA CFD Drag Prediction Workshop (DPW-IV), the NASA Common Research Model (CRM) has become one of the reference configurations for the validation of CFD methods applied to transonic transport aircraft~\cite{vassberg2008}. It combines a contemporary supercritical wing, with an aspect ratio of 9 and a quarter-chord sweep angle of \SI{35}{\degree}, with a fuselage representative of a wide-body airliner, and it was designed to cruise at $M_\infty = 0.85$ with a lift coefficient $C_L = 0.5$. The geometry is publicly available and, over the years, the model has been tested in several transonic wind tunnels and simulated with a wide variety of solvers, grids and turbulence models, many of them within the DPW series. The resulting body of experimental and numerical data makes the CRM a natural starting point for this work, which moves from the steady cruise condition to the unsteady buffet regime.

Before turning to high-fidelity simulations of buffet, however, it is necessary to verify that the numerical chain as a whole, from the computational grid to the physical model and the flow solver, gives reliable results in a simpler situation. The design cruise point is well suited to this purpose, since the flow is steady and largely attached and reference data are abundant. This chapter presents such a validation, carried out by means of steady Reynolds-Averaged Navier-Stokes (RANS) simulations of the CRM in the wing-body-tail configuration (WBT0) at $M_\infty = 0.85$ and $Re = 5\times10^{6}$, based on the mean aerodynamic chord.

The work was carried out in two stages. In the first one, the computational grids generated with Ansys Fluent Meshing were assessed with the commercial solver Ansys Fluent. Relying on a well-established code at this stage made it possible to evaluate the grids and the meshing strategy on their own, by comparing the results with the experimental measurements of Rivers et al.~\cite{rivers2015} and with the numerical results of Du and Agarwal~\cite{du2017}, which were obtained with the same software. After a mesh sensitivity analysis, the reference grid was used to compare the $k$--$\omega$ and Spalart--Allmaras turbulence models and to perform an angle-of-attack sweep from \SI{0}{\degree} to \SI{3}{\degree}. Higher incidences were not considered, since the growth of shock-induced separation progressively drives the flow towards buffet, a regime that a steady RANS approach cannot describe.

In the second stage, the same grid was used to run Spalart--Allmaras simulations with STREAmS, the in-house solver developed at Sapienza University of Rome that is also employed for the high-fidelity simulations presented later in this thesis. Since the grid and the turbulence model are the same as in the Fluent computations, the differences between the two solutions can be attributed to the solver itself. The STREAmS results are compared with both the Fluent solution and the experimental data.

The chapter follows the same order: after the description of the test case and of the numerical setup, the Fluent results are presented and discussed, and the validation of STREAmS concludes the chapter.

\section{Test case description}
\label{sec:val_testcase}

\subsection{Geometry and reference quantities}
\label{sec:val_geometry}

The geometry considered in this work is the wing/body/horizontal tail configuration of the CRM with the stabiliser set at \SI{0}{\degree} incidence (WBT0), in the full-scale version released for DPW-IV~\cite{vassberg2008}. The models tested in wind tunnels are scaled replicas of this geometry: the one used in the ETW and NTF campaigns has a mean aerodynamic chord of \SI{7.45}{in}, which corresponds to a scale factor of about 2.7\%~\cite{rivers2015}. Since the flow condition is defined through non-dimensional parameters, computations on the full-scale geometry and measurements on the scaled models refer to the same aerodynamic condition, as discussed in Section~\ref{sec:val_flow}.

Since both the configuration and the flow at zero sideslip are symmetric with respect to the vertical plane containing the fuselage axis, only half of the aircraft is simulated and a symmetry condition is imposed on that plane. Aerodynamic coefficients are therefore computed with half of the reference area, so that they correspond to those of the complete aircraft. The reference quantities are summarised in Table~\ref{tab:val_refquantities}.
% TODO: se si calcola C_M, aggiungere: "the pitching moment is evaluated about the moment reference centre of the DPW-IV geometry" + riga in tabella

\begin{table}[htbp]
  \centering
  \caption{Reference quantities of the full-scale CRM.}
  \label{tab:val_refquantities}
  \begin{tabular}{lcc}
    \hline
    Quantity & Value & SI units \\
    \hline
    Reference area, $S_\mathrm{ref}$ & \SI{594720}{in^2} & \SI{383.69}{m^2} \\
    Half-model reference area, $S_\mathrm{ref}/2$ & \SI{297360}{in^2} & \SI{191.845}{m^2} \\
    Mean aerodynamic chord, $c_\mathrm{ref}$ & \SI{275.8}{in} & \SI{7.005}{m} \\
    Wing span, $b$ & \SI{2313.4}{in} & \SI{58.76}{m} \\
    Aspect ratio, $AR$ & \multicolumn{2}{c}{9.0} \\
    \hline
  \end{tabular}
\end{table}

\subsection{Flow conditions}
\label{sec:val_flow}

The simulations are performed at the reference condition of the CRM, $M_\infty = 0.85$, $C_L = 0.5$ and $Re = 5\times10^{6}$, based on $c_\mathrm{ref}$. This is the condition defined for the main test case of DPW-IV, which combines the design Mach number and lift coefficient of the aircraft with a Reynolds number representative of wind tunnel testing~\cite{vassberg2010}, lower than the one typical of flight. The same condition was retained in the following editions of the workshop and has been adopted both in the experimental campaigns on the CRM~\cite{rivers2015} and in most of the numerical studies of this configuration, so that it has become the baseline for any validation on the CRM.

Since the DPW-IV geometry is defined at full scale, with $c_\mathrm{ref} = \SI{7.005}{m}$, the freestream state that reproduces these conditions has to be derived from the similarity parameters. The DPW-IV specification also fixes the reference temperature at \ang{100}F; accordingly, the freestream static temperature is set to $T_\infty = \SI{310.83}{K}$ and the dynamic viscosity is computed with Sutherland's law,
\begin{equation}
  \mu_\infty = \mu_0 \left( \frac{T_\infty}{T_0} \right)^{3/2} \frac{T_0 + S}{T_\infty + S},
  \label{eq:val_sutherland}
\end{equation}
with $\mu_0 = \SI{1.716e-5}{kg/(m.s)}$, $T_0 = \SI{273.11}{K}$ and $S = \SI{110.56}{K}$. The freestream velocity follows from the Mach number, treating air as a calorically perfect gas with $\gamma = 1.4$, while the density is obtained from the definition of the Reynolds number and the pressure from the equation of state:
\begin{equation}
  U_\infty = M_\infty \sqrt{\gamma R T_\infty}, \qquad
  \rho_\infty = \frac{Re\, \mu_\infty}{U_\infty\, c_\mathrm{ref}}, \qquad
  p_\infty = \rho_\infty R T_\infty .
  \label{eq:val_freestream}
\end{equation}
The resulting values are listed in Table~\ref{tab:val_flowconditions}. The low density and pressure, far from those encountered by the aircraft at cruise altitude, are a direct consequence of imposing a wind tunnel Reynolds number on a full-scale geometry.

\begin{table}[htbp]
  \centering
  \caption{Freestream conditions used in the simulations.}
  \label{tab:val_flowconditions}
  \begin{tabular}{lc}
    \hline
    Quantity & Value \\
    \hline
    Mach number, $M_\infty$ & 0.85 \\
    Reynolds number, $Re$ (based on $c_\mathrm{ref}$) & $5\times10^{6}$ \\
    Static temperature, $T_\infty$ & \SI{310.83}{K} \\
    Dynamic viscosity, $\mu_\infty$ & \SI{1.897e-5}{kg/(m.s)} \\
    Velocity, $U_\infty$ & \SI{300.3}{m/s} \\
    Density, $\rho_\infty$ & \SI{0.0451}{kg/m^3} \\
    Static pressure, $p_\infty$ & \SI{4022}{Pa} \\
    \hline
  \end{tabular}
\end{table}

The reference cruise condition is computed at an angle of attack $\alpha = \SI{2.38}{\degree}$, for which the Fluent solution gives a lift coefficient very close to the design value $C_L = 0.5$. The same incidence is used for all the comparisons at cruise conditions, including those between the two turbulence models and between Fluent and STREAmS. The boundary layer is assumed to be fully turbulent: in the experiments transition was fixed by trip dots placed at 10\% of the local chord~\cite{rivers2015}, so that the laminar region is limited to a small portion close to the leading edge.

\subsection{Reference data}
\label{sec:val_refdata}

The results are compared with two kinds of reference data: experimental measurements and numerical results from the literature. The sources used and the quantities taken from each of them are summarised in Table~\ref{tab:val_refdata}.

The experimental data are taken from the work of Rivers et al.~\cite{rivers2015}, who compared the measurements obtained on the same CRM model in three facilities, the NASA Langley National Transonic Facility (NTF), the NASA Ames \SI{11}{ft} Transonic Wind Tunnel and the European Transonic Wind Tunnel (ETW), at different conditions, including the DPW4 main test case, which is the one investigated here. 

The main comparison concerns the wing pressure distributions, which describe in detail the flow at the cruise point. Among the measurements available at $M_\infty = 0.85$ and $Re = 5\times10^{6}$, the ETW data point with the lift coefficient closest to the design value was selected (run 153, point 345), corresponding to $\alpha = \SI{2.88}{\degree}$ and $C_L = 0.498$. Pressure distributions are reported for this point at six spanwise stations, $\eta = 0.201$, $0.283$, $0.502$, $0.603$, $0.846$ and $0.950$, where $\eta$ is the spanwise coordinate normalised by the semi-span, and all of them are used in the comparison; the measurement uncertainty on $C_p$ is estimated at $\pm 0.0026$~\cite{rivers2015}.

The same work reports the lift, drag and pitching moment coefficients measured in the three facilities as functions of the angle of attack. These data are used, first, to assess the integral coefficients at the cruise point and, secondly, as a reference for the angle-of-attack sweep. The three data sets do not coincide exactly: at $M_\infty = 0.85$, for instance, the drag measured in the NTF is 16 counts lower than that measured in the ETW~\cite{rivers2015}. The spread among the facilities is therefore taken as a practical measure of the experimental uncertainty.

The experimental lift coefficient is reached at an angle of attack about half a degree higher than in the computations. For this reason the pressure distributions are compared at matched lift coefficient rather than at matched incidence, which is also the criterion followed in~\cite{rivers2015} to compare data from different facilities. The difference in incidence reflects effects that are present in the experiment but not in the simulations, such as the deformation of the wing under load, which was measured during the tests, and the interference of the model support, for which the published data were not corrected~\cite{rivers2015}; this point is discussed further in Section~\ref{sec:val_cruise}.

Among the many numerical studies of the CRM, three are used as references. The first is DPW-IV itself, whose test cases correspond to the configuration and flow condition considered here. The statistical analysis of the solutions submitted to the workshop~\cite{morrison2010} provides the median and the scatter, over about twenty RANS solutions obtained with different codes, of the angle of attack at $C_L = 0.5$, of the drag coefficient, split into its pressure and friction contributions, and of the pitching moment; the lift curves and drag polars computed for the angle-of-attack sweep of the workshop~\cite{vassberg2010} complete this set. The second reference is the work of Du and Agarwal~\cite{du2017}, who computed the same configuration with Ansys Fluent using both the Spalart--Allmaras and the SST $k$--$\omega$ models on a structured grid of about 11 million nodes; since it relies on the same solver, it allows a direct comparison of grids and solver settings.
% PROVVISORIO: frase su Sclafani et al. (da togliere, insieme alla riga in tabella, se il paper non viene usato)
Finally, Sclafani et al.~\cite{sclafani2010} compared the two turbulence models on the same configuration with the CFL3D and OVERFLOW solvers, and their results are used to interpret the differences between the SST $k$--$\omega$ and Spalart--Allmaras solutions.

\begin{table}[htbp]
  \centering
  \caption{Reference data used in this chapter.}
  \label{tab:val_refdata}
  \small
  \begin{tabular}{>{\raggedright\arraybackslash}p{3.2cm}>{\raggedright\arraybackslash}p{2.5cm}>{\raggedright\arraybackslash}p{4.0cm}>{\raggedright\arraybackslash}p{2.0cm}}
    \hline
    Source & Type & Data used & Section \\
    \hline
    Rivers et al.~\cite{rivers2015}, ETW run 153, point 345 & Experimental & $C_p$ at six spanwise stations ($\alpha = \SI{2.88}{\degree}$, $C_L = 0.498$) & \ref{sec:val_cruise} \\[2pt]
    Rivers et al.~\cite{rivers2015}, NTF, Ames and ETW & Experimental & $C_L$, $C_D$ and $C_M$ versus $\alpha$ & \ref{sec:val_cruise}, \ref{sec:val_alpha} \\[2pt]
    DPW-IV~\cite{vassberg2010,morrison2010} & Numerical, workshop ensemble & $\alpha$ at $C_L = 0.5$, $C_D$ (pressure and friction), $C_M$; lift curves and polars & \ref{sec:val_mesh}, \ref{sec:val_cruise}, \ref{sec:val_alpha} \\[2pt]
    Du and Agarwal~\cite{du2017} & Numerical, Fluent (SA, SST) & $C_L$ and $C_D$ versus $\alpha$, drag polar & \ref{sec:val_cruise}, \ref{sec:val_turb}, \ref{sec:val_alpha} \\[2pt]
    % PROVVISORIO: riga Sclafani et al.
    Sclafani et al.~\cite{sclafani2010} & Numerical, CFL3D and OVERFLOW (SA, SST) & Effect of the turbulence model on shock position, $C_D$, $\alpha$ and $C_p$ & \ref{sec:val_turb} \\
    \hline
  \end{tabular}
\end{table}


\section{Grid generation strategy}
\label{sec:val_grid}
% TODO: Fluent Meshing, tipo di elementi, inflation layer (y+ ~ 1),
%       growth rate, zone di raffinamento (LE/TE, urto, side-of-body),
%       dimensione del dominio e farfield.

\section{Numerical setup}
\label{sec:val_setup}
% TODO: solver density/pressure-based, schemi, k-omega (SST?),
%       condizioni al contorno, inizializzazione, criteri di convergenza.
% TODO: giustificazione di alpha fissato invece di CL fissato.

\section{Mesh sensitivity analysis}
\label{sec:val_mesh}
% TODO: tabella famiglia di griglie (n. celle, y+, ...),
%       CL/CD/CM vs N^{-2/3}, eventuale Richardson/GCI, mesh scelta.

\section{Validation at cruise conditions}
\label{sec:val_cruise}
\subsection{Integral coefficients}
% TODO: CL, CD (pressione/attrito), CM vs paper numerico.
\subsection{Pressure distributions}
% TODO: Cp alle stazioni eta vs sperimentale e numerico; posizione urto.

\section{Turbulence model sensitivity}
\label{sec:val_turb}
% TODO: k-omega vs SA alla condizione di crociera.

\section{Angle-of-attack sweep}
\label{sec:val_alpha}
% TODO: alpha = 0-3 deg; CL-alpha, polare, CM-alpha vs sperimentale/numerico;
%       evoluzione dell'urto; limite a 3 deg per vicinanza del buffet.

\section{Validation of STREAmS}
\label{sec:val_streams}
\subsection{Solver and setup}
% TODO: richiamo al capitolo dei metodi, trasferimento della mesh, setup SA.
\subsection{Comparison with Fluent and experimental data}
% TODO: coefficienti integrali + Cp, STREAmS vs Fluent-SA vs sperimentale.

\section{Summary and implications for the unsteady analysis}
\label{sec:val_summary}
% TODO: sintesi dell'accordo, limiti, cosa si porta nella fase WMLES.
